
\subsection*{\SGARefTarget{sga7.sga7i.evii.n3-7}{3.7}. Various compatibilities}

In this section we examine the behaviour, for variable $P$ and $Q$,
of the equivalence of cofibred categories in \SGARefLink{sga7.sga7i.evii.n3-6-4}{3.6.4}.  The variance of
the cofibred category of biextensions of $(P,Q)$ by a variable $G$
with respect to variable $P,Q$ is defined in principle in \SGARefLink{sga7.sga7i.evii.n2-4}{2.4}; it may
also be expressed by saying that one has a cofibred functor (\SGARefLink{sga7.sga7i.evii.n2-6-1}{2.6.1}).
Notice that the variance in $P,Q$ is \underline{contravariant}.  The
variance of the cofibred category
$\underline\Psi_{L_\bullet(P,Q)}$ with respect to variable $P,Q$ is
also clear, since
$L_\bullet(P,Q)=L_\bullet(P)\otimes L_\bullet(Q)$ depends
functorially on the pair $(P,Q)$, while $\underline\Psi_{L_\bullet}$
depends contravariantly on $L_\bullet$ (\SGARefLink{sga7.sga7i.evii.n3-4}{3.4}).  With this understood,
a first compatibility is expressed, for every triple of morphisms
\begin{equation}\tag{\SGARefTarget{sga7.sga7i.evii.n3-7-1}{3.7.1}}
 P'\longrightarrow P,\qquad
 Q'\longrightarrow Q,\qquad
 G\longrightarrow G',
\end{equation}
by the fact that the diagram of functors
\begin{equation}\tag{\SGARefTarget{sga7.sga7i.evii.n3-7-2}{3.7.2}}
\begin{tikzcd}[column sep=huge,row sep=large]
 \underline\Psi_{L_\bullet(P,Q)}(G)
   \arrow[r]\arrow[d]
 &\underline\Psi_{L_\bullet(P',Q')}(G')
   \arrow[d]\\
 \operatorname{BIEXT}(P,Q;G)
   \arrow[r]
 &\operatorname{BIEXT}(P',Q';G')
\end{tikzcd}
\end{equation}
is commutative up to a canonical isomorphism.  We shall admit this
compatibility: making the commutativity isomorphism explicit could
only be as obvious as it is tedious (as, indeed, could the complete
explication it presupposes of the equivalence of cofibred categories
in \SGARefLink{sga7.sga7i.evii.n3-6-4}{3.6.4}).

\SGARefTarget{sga7.sga7i.evii.n3-7-3}{3.7.3}. It follows in particular from the preceding compatibility
that the isomorphism (\SGARefLink{sga7.sga7i.evii.n3-6-5-1}{3.6.5.1}) is functorial not only in $G$, but also
in $P$ and $Q$.  One may state, similarly and more trivially, the
analogous functoriality of the isomorphism defined through (\SGARefLink{sga7.sga7i.evii.n3-1-8}{3.1.8}),
just as (\SGARefLink{sga7.sga7i.evii.n3-6-5-1}{3.6.5.1}) was defined through (\SGARefLink{sga7.sga7i.evii.n3-2-4}{3.2.4}):
\begin{equation}\tag{\SGARefTarget{sga7.sga7i.evii.n3-7-4}{3.7.4}}
 \Biext^0(P,Q;G)
 \xrightarrow{\sim}
 \Ext^0(P\mathbin{\overset{\mathbb L}{\otimes}}Q,G)
 \simeq\Hom(P\otimes Q,G).
\end{equation}
But such an isomorphism was already obtained in (\SGARefLink{sga7.sga7i.evii.n2-5-4-1}{2.5.4.1}), and was
observed in \SGARefLink{sga7.sga7i.evii.n2-6}{2.6} to be functorial in its three arguments.  It is
immediate to verify that the two isomorphisms (\SGARefLink{sga7.sga7i.evii.n3-7-4}{3.7.4}) and (\SGARefLink{sga7.sga7i.evii.n2-5-4-1}{2.5.4.1})
are equal.

\SGARefTarget{sga7.sga7i.evii.n3-7-5}{3.7.5}. The value of the isomorphism in \SGARefLink{sga7.sga7i.evii.n3-6-5}{3.6.5} is that it expresses the
remarkable group of geometric origin $\Biext^1(P,Q;G)$ in terms of
the very simple homological-algebra invariant
$\Ext^1(P\mathbin{\overset{\mathbb L}{\otimes}}Q,G)$.  The latter
belongs to the sequence of invariants
$\Ext^i(P\mathbin{\overset{\mathbb L}{\otimes}}Q,G)$, to which the
usual formalism of homological algebra applies.  In particular, the
sequence of these $\Ext^i$ is a \underline{cohomological functor} in
the sense of [3] (or a \underline{connected sequence of functors} in
the terminology of [5]) in each of the three arguments $P,Q,G$.
Thus an exact sequence
\begin{equation}\tag{\SGARefTarget{sga7.sga7i.evii.n3-7-5-1}{3.7.5.1}}
 0\longrightarrow G'\longrightarrow G\longrightarrow G''
 \longrightarrow0
\end{equation}
(respectively,
\begin{equation}\tag{\SGARefTarget{sga7.sga7i.evii.n3-7-5-2}{3.7.5.2}}
 0\longrightarrow P'\longrightarrow P\longrightarrow P''
 \longrightarrow0,\qquad\text{or}\qquad
 0\longrightarrow Q'\longrightarrow Q\longrightarrow Q''
 \longrightarrow0)
\end{equation}
gives rise to an infinite exact sequence of the $\Ext^i$.  Its first
six terms, which involve only $\Ext^0$ and $\Ext^1$, may be
interpreted in terms of biextensions by \SGARefLink{sga7.sga7i.evii.n3-6-5}{3.6.5} and \SGARefLink{sga7.sga7i.evii.n3-7-4}{3.7.4}.  We have
already indicated in \SGARefLink{sga7.sga7i.evii.n3-3-6}{3.3.6} a geometric interpretation of the exact
sequence deduced from (\SGARefLink{sga7.sga7i.evii.n3-7-5-1}{3.7.5.1}), and especially of the boundary map
\begin{equation}\tag{\SGARefTarget{sga7.sga7i.evii.n3-7-5-3}{3.7.5.3}}
 \partial:\Biext^0(P,Q;G'')\longrightarrow
 \Biext^1(P,Q;G'),
\end{equation}
within the framework of a \underline{left exact cofibred category}.
One may seek an analogous interpretation of the six-term exact
sequence associated with either short exact sequence in (\SGARefLink{sga7.sga7i.evii.n3-7-5-2}{3.7.5.2}),
say the first; equivalently, in view of \SGARefLink{sga7.sga7i.evii.n3-7-3}{3.7.3}, one may seek a
geometric interpretation of the boundary homomorphism
\begin{equation}\tag{\SGARefTarget{sga7.sga7i.evii.n3-7-5-4}{3.7.5.4}}
 \partial:\Biext^0(P',Q;G)\longrightarrow
 \Biext^1(P'',Q;G).
\end{equation}
Indeed, we have the following result.

\underline{Proposition} \SGARefTarget{sga7.sga7i.evii.n3-7-6}{3.7.6}. \underline{The cofibred category over}
\[
 \operatorname{GRAB}(\underline T)^{\mathrm{op}}
 \times\operatorname{GRAB}(\underline T)^{\mathrm{op}}
 \times\operatorname{GRAB}(\underline T)
\]
\underline{of biextensions of $(P,Q)$ by $G$, with $P,Q,G$
variable} (cf. (\SGARefLink{sga7.sga7i.evii.n2-6-1}{2.6.1})), \underline{is left exact not only in the
argument $G$, but also in each of the arguments $P$ and $Q$.
Moreover, the boundary homomorphism (\SGARefLink{sga7.sga7i.evii.n3-7-5-4}{3.7.5.4}) deduced from the exact
sequence of the $\Ext^i$ associated with (\SGARefLink{sga7.sga7i.evii.n3-7-5-2}{3.7.5.2}), by means of the
isomorphisms \SGARefLink{sga7.sga7i.evii.n3-6-5}{3.6.5} and \SGARefLink{sga7.sga7i.evii.n3-7-4}{3.7.4}, is the boundary homomorphism in the
theory of left exact cofibred categories over}
$\underline C=\operatorname{GRAB}(\underline T)^{\mathrm{op}}$
\underline{recalled in \SGARefLink{sga7.sga7i.evii.n3-3-6}{3.3.6}, applied to the cofibred category of
biextensions of $(P,Q)$ by $G$, where $Q$ and $G$ are fixed and $P$
varies} (subject to a possible sign error in \SGARefLink{sga7.sga7i.evii.n3-4-9}{3.4.9}; cf. \SGARefLink{sga7.sga7i.evii.n3-4-9-2}{3.4.9.2}).

We sketch the proof of \SGARefLink{sga7.sga7i.evii.n3-7-6}{3.7.6}.  By \SGARefLink{sga7.sga7i.evii.n3-6-4}{3.6.4} and compatibility (\SGARefLink{sga7.sga7i.evii.n3-7-2}{3.7.2}),
we interpret biextensions of $(P,Q)$ by $G$ as objects of
$\underline\Psi_{L_\bullet(P,Q)}(G)$.  Consider the homomorphism
\[
 L_\bullet(P,Q)
 =L_\bullet(P)\otimes L_\bullet(Q)
 \longrightarrow P\otimes L_\bullet(Q)
\]
defined by the augmentation $L_\bullet(P)\longrightarrow P$.
It follows from \SGARefLink{sga7.sga7i.evii.n3-5-3}{3.5.3} and the flatness of $L_\bullet(Q)$ that this
homomorphism induces isomorphisms on $H_0$ and $H_1$, and hence by
\SGARefLink{sga7.sga7i.evii.n3-4-7}{3.4.7} an equivalence
\[
 \underline\Psi_{P\otimes L_\bullet(Q)}
 \longrightarrow
 \underline\Psi_{L_\bullet(P,Q)}.
\]
We may therefore interpret biextensions of $(P,Q)$ by $G$ as objects
of $\underline\Psi_{P\otimes L_\bullet(Q)}(G)$.  The exact sequence
(\SGARefLink{sga7.sga7i.evii.n3-7-5-2}{3.7.5.2}) now gives an exact sequence of chain complexes
\begin{equation}\tag{\SGARefTarget{sga7.sga7i.evii.n3-7-6-1}{3.7.6.1}}
 0\longrightarrow P'\otimes L_\bullet(Q)
 \longrightarrow P\otimes L_\bullet(Q)
 \longrightarrow P''\otimes L_\bullet(Q)
 \longrightarrow0.
\end{equation}
The asserted left exactness follows from \SGARefLink{sga7.sga7i.evii.n3-4-8}{3.4.8} applied to this exact
sequence of complexes.  (A direct verification would in any event
be immediate.)  Moreover, the homomorphism (\SGARefLink{sga7.sga7i.evii.n3-7-5-4}{3.7.5.4}) is precisely the
boundary homomorphism $\Ext^0\longrightarrow\Ext^1$ deduced from this
exact sequence by taking the $\Ext^i$ with values in $G$.  The last
assertion of \SGARefLink{sga7.sga7i.evii.n3-7-6}{3.7.6} is consequently a special case of \SGARefLink{sga7.sga7i.evii.n3-4-9}{3.4.9}.
